\begin{table}[htbp]
\centering
\footnotesize
\caption{Selected prior work and the space occupied by this study. The table is a
focused comparison, not a systematic review. It shows that the individual
ingredients of the Jakarta index have precedent; the contribution claimed here is
the attribution and uncertainty audit, not the combination itself.}
\label{tab:literature}
\begin{threeparttable}
\begin{tabularx}{\textwidth}{@{}>{\raggedright\arraybackslash}p{2.45cm}
>{\raggedright\arraybackslash}p{2.3cm}
>{\raggedright\arraybackslash}p{2.5cm}
>{\raggedright\arraybackslash}p{2.0cm}
>{\raggedright\arraybackslash}p{2.1cm}
>{\raggedright\arraybackslash}X@{}}
\toprule
Study & Context / unit & Construct & Accessibility & Uncertainty & Relation to this study \\
\midrule
Oke (1982); Voogt \& Oke (2003) & Urban climate; conceptual & Surface energy balance; thermal remote sensing & None & Not addressed & Conceptual basis for LST; no prioritisation \\
Weng et al.\ (2004); Estoque et al.\ (2017) & Cities; pixel or grid & LST--vegetation relationship & None & Limited & Supports M1--M2 layers; no demographic or network term \\
Reid et al.\ (2009); Wolf \& McGregor (2013) & US; London; areal units & Heat vulnerability index & None or facility proximity & Limited & Composite vulnerability precedent; no rank attribution \\
Sidiqui et al.\ (2022) & GIS-based urban case & Heat hazard, exposure and sensitivity & Facility/park access & Scenario-limited & Closest composite precedent; not Jakarta, not attribution-focused \\
Kim \& Kwak (2026) & Older adults; walking network & Cooling-weighted park accessibility (G2SFCA) & OSM walking network & Demand and network sensitivity & Close overlap in access and age; no Jakarta case or rank decomposition \\
Fung et al.\ (2024) & Urban heat mitigation & Social-vulnerability prioritisation & None & Scenario comparison & Prioritisation precedent; not park-network uncertainty \\
Sobhaninia et al.\ (2025) & Systematic review & Green-infrastructure siting methods & Various & Review-level & Positions indicator choices; no Jakarta application \\
Rizki et al.\ (2024) & East Jakarta; LST/NDVI & Green space and surface heat island & None & Limited & Jakarta evidence; no demographic or network attribution \\
Varquez et al.\ (2020) & Jakarta; mortality projection & Elderly heat-related mortality & None & Climate/population projection & Health-impact estimand, not a screening index \\
This study & Five mainland Jakarta cities; 500 m grid & Attribution and uncertainty audit of a heat--population--park screening index & OSM candidate-park network proxy & 52 scenarios; Shapley; LISA conventions & Conditional contribution: attribution, negative findings, reproducibility \\
\bottomrule
\end{tabularx}
\begin{tablenotes}[flushleft]\footnotesize
\item G2SFCA: generalised two-step floating catchment area. The table summarises what each study demonstrates, not every method used. The absence of a study from this table is not evidence that no similar work exists.
\end{tablenotes}
\end{threeparttable}
\end{table}
